\section{An all-index large-prime calculation and a weight-dependent family}
\label{sec:forced-hankel-depth}

The local calculation below tests the proposed transfer from weight five
using the actual monic polynomial. It involves no finite-index experiments.
Write $P_{s,K}=\Delta_{s,K}/L_{s,K}$ and let $d_{s,6}(K)$ be the least
positive integer clearing its coefficients. Thus
\[
 v_p(d_{s,6}(K))=-v_p^G(P_{s,K}).
\]

\begin{proposition}\label{prop:forced-large-prime-depth}
Fix odd $s\ge9$. For every sufficiently large $K=40n$ and every prime
$2K/3<p<K$, put $m=K-p+1$ and $h=K-N$, where $N=3n$.
Then
\[
 v_p\bigl([X^{h-m}]P_{s,K}\bigr)=-sm,
 \qquad p^{sm}\mid d_{s,6}(K).                         \tag{D1}
\]
Consequently
\[
 \liminf_{n\to\infty}\frac{\log d_{s,6}(40n)}{(40n)^2}
 \ge \frac{s}{18}.                                    \tag{D2}
\]
\end{proposition}
\begin{proof}
All reductions in this proof are over $\Z_{(p)}$ and its residue field.
For sufficiently large $K$, we have $p>s+1$ and $p^2>2K$.
Set $I=\{N+1,\ldots,K\}$,
\[
 E(t)=\prod_{i\in I}(t+i^2),\quad e_i=E'(-i^2),\quad
 \ell_i(t)=\frac{E(t)}{(t+i^2)e_i}.
\]
The Lagrange basis diagonalizes the coefficient of $X$ in the Gram
matrix. Its diagonal entries are
$b_i=i^{s-1}D_N(-i^2)^5/e_i$, all nonzero.
After division by this diagonal matrix, followed by conjugation by
$\operatorname{diag}(e_i)$, the monic polynomial becomes
\[
 P_{s,K}(X)=\det(XI+A),
\]
where
\[
 A_{ij}=\frac{\mu_s(Q_{ij})}
 {e_i i^{s-1}D_N(-i^2)^5}
 +\delta_{ij}\left(
 \frac{1/(2i)-1/(s-1)}{i^{s-1}}-H_i^{(s)}\right).       \tag{D3}
\]
Here $\mu_s$ denotes the polynomial-moment functional and
\[
 Q_{ij}=\frac{D_N(t)^5E(t)}{(t+i^2)(t+j^2)}\quad(i\ne j),
\]
\[
 Q_{ii}=\frac{D_N(t)^5 E(t)/(t+i^2)-D_N(-i^2)^5e_i}
 {t+i^2}.
\]
Both formulas define polynomials in $\Z[t]$: in the second formula
the numerator vanishes at $-i^2$, so division by the monic factor is
exact. Their degrees are at most $K+4N-2$.

We have $\mu_s(Q_{ij})\in\Z_{(p)}$. Indeed
\[
 2(K+4N-2)+2 < (s-1)(p-1)
\]
for the stated range and sufficiently large $K$. The factorial
$(s-1)!$ is a $p$-unit. To check the cancellation directly,
put $L=2e+2$. Von Staudt--Clausen makes $B_L$ integral unless
$p-1\mid L$. In the latter case write $L=k(p-1)$; the inequality
gives $1\le k\le s-2$, and the numerator factor $L+k=kp$ cancels
its possible denominator $p$.

Since $K<3p/2$, a nonzero square class among $1,\ldots,K$ has at most
three representatives. Each nonzero difference $i^2-j^2$ in the same
class has valuation one, because $p^2>2K$. Hence
$v_p(e_i)\le2$ for $p\nmid i$. Also
$v_p(D_N(-i^2))\le1$, since $2N<p$ permits at most one matching
root in $1,\ldots,N$. For the only multiple of $p$ in $I$, namely
$i=p$, both $e_i$ and $D_N(-i^2)$ are units.
It follows that the first term of (D3) has valuation at least $-7$
when $i\ne p$, and at least $-(s-1)$ when $i=p$.
Thus multiplication by $p^s$ makes every such entry zero modulo $p$.

For $i<p$ the remaining diagonal term is integral. For $i>p$ its
product with $p^s$ is $-1$ modulo $p$: the only nonintegral term of
$H_i^{(s)}$ is $p^{-s}$, because $i<2p$.
At $i=p$ the boundary term must be retained; it gives
\[
 p^s\left(
 \frac{1/(2p)-1/(s-1)}{p^{s-1}}-p^{-s}
 \right)=-\frac12-\frac{p}{s-1}\equiv-\frac12\pmod p.
\]
We have therefore proved the exact polynomial congruence
\[
 p^{sh}P_{s,K}(X/p^s)
 \equiv X^{h-m}(X-\tfrac12)(X-1)^{m-1}\pmod p.         \tag{D4}
\]
The coefficient of $X^{h-m}$ on the right is $(-1)^m/2$, a unit.
On the left it is $p^{sm}[X^{h-m}]P_{s,K}$. This proves (D1).
As an independent way of checking this coefficient, in the elementary
symmetric expansion of the reduced diagonal matrix there is exactly
one nonzero term of order $m$: it selects all $m$ nonzero entries.

Finally the prime number theorem and partial summation give
\[
 \sum_{2K/3<p<K}(K-p+1)\log p
 =K^2\int_{2/3}^1(1-u)\,du+o(K^2)
 =\frac{K^2}{18}+o(K^2).
\]
Summing (D1) proves (D2).
\end{proof}

\paragraph{The same calculation at weight seven.}
There is a useful weaker conclusion at the primary target weight.
For $2K/3<p<K$, the matrix $p^7A$ is integral, but its polynomial
part need not vanish on rows where $p\mid D_N(-i^2)$.
Call these rows bad. All other rows reduce to their diagonal pole
term. The bad active indices $i>p$ are exactly
$p+v$, $1\le v\le\min(N,K-p)$. The index $i=p$ is good.
Thus at least
\[
 g=\max(0,K-p-N)+1
\]
good diagonal entries are nonzero. Order the bad rows before the
good rows. The reduced matrix is block triangular, and its
characteristic polynomial contains a degree-$g$ factor with unit
constant coefficient from these good active entries. The remaining
factor is monic. Its lowest nonzero coefficient shows that some
coefficient $[X^k]p^{7h}P_{7,K}(X/p^7)$ is a unit with $k\le h-g$.
Consequently
\[
 p^{7g}\mid d_{7,6}(K),\qquad
 \liminf\frac{\log d_{7,6}(K)}{K^2}
 \ge\frac72\left(\frac13-\frac3{40}\right)^2
 =\frac{6727}{28800}.                                  \tag{D7}
\]
Here polynomial-moment integrality follows from
$2(K+4N-2)+2<6(p-1)$. This lower bound, like (D2), measures a
cost that a successful normalized decay estimate must include.

\paragraph{Consequence for a fixed-power transfer.}
This is a lower bound on a primitive denominator, rather than on
an arbitrary clearing factor. It therefore survives changes of basis.
For completeness a deliberately coarse upper bound at weight five is
\[
 \limsup\frac{\log d_{5,6}(K)}{K^2}\le\frac{139}{4}.
                                                               \tag{D5}
\]
To verify the constant without numerical experiments, Proposition~
\ref{prop:coarse-local-hankel} shows that
$C_KF_{5,K}\in\Z[X]$ for
\[
 C_K=\operatorname{lcm}(1,\ldots,5K)^{6h}24^h,
 \qquad \frac{\log C_K}{K^2}\longrightarrow\frac{111}{4}.
\]
The monic coefficient denominator divides the absolute value of the
leading coefficient of this integer polynomial. Put $a=3/40$ and
$l=37/40$. The leading coefficient satisfies
$\log|L_{5,K}|\le(10Nh+4h)\log K$, since
$0<j^2-v^2\le K^2$ for $N<j\le K$, $1\le v\le N$.
Stirling's formula for $S_K$, with cancellation of the
$K^2\log K$ terms, therefore gives
\[
 \limsup\frac{\log|S_KL_{5,K}|}{K^2}
 \le-2l-12la\log a+12la-2l^2\log(2l)+3l^2.
\]
Discard the negative terms and use $\log(40/3)<3$ to bound this above by
$36a+12a+3=33/5<7$. This proves (D5).
Combining (D2) and (D5) gives
\[
 \liminf\frac{\log(d_{s,6}(K)/d_{5,6}(K))}{K^2}
 \ge\frac{s}{18}-\frac{139}{4}.                         \tag{D6}
\]
In particular the fixed sixth-power family cannot have relative cost
below one fixed finite margin for every odd weight. This conclusion
does not rule out the weight-seven transfer or the irrationality of
any zeta value. It rules out that specific all-weight estimate for
these fixed parameters.

\paragraph{An explicit family with weight-dependent power.}
The calculation also supplies a concrete design choice. Replace
$D_N^6$ by $D_N^{s+1}$, retaining $K=40n$, $N=3n$ and $h=37n$:
\[
 \widetilde\Delta_{s,K}(X)=\det\left[
 \mu_{s,X}\left(\frac{D_N(t)^{s+1}t^{i+j}}{D_K(t)}\right)
 \right]_{0\le i,j<h}.
\]
This is the existing weight-five family when $s=5$. Its leading
coefficient and polynomial quotient degree are explicitly
\[
 \widetilde L_{s,K}=(-1)^{h(h-1)/2}
 \prod_{j=N+1}^K j^{s-1}D_N(-j^2)^s,
 \qquad \widetilde E_s=K+(s-1)N-2.
\]
The same positive measure representation proves
$\widetilde\Delta_{s,K}(\zeta(s))>0$ for every odd $s\ge5$.
In (D3) the factor $D_N(-i^2)^5$ becomes $D_N(-i^2)^s$.
Thus the domination used in (D4) is no longer automatic. This
constructs a family where numerator multiplicity can compensate
for the increasing harmonic pole order. It does not prove the
compensation or the required negative total rate: both the new
local prime allocation and the logarithmic-energy estimate must
be established in this normalization.

\paragraph{Use of recent prime-distribution results.}
The OpenAI preprint of 30 September 2026 states a zero-free half-plane
for Dirichlet and specified Hecke $L$-functions; see
\href{https://github.com/openai/math/blob/main/preprints/The-Quasi-Riemann-Hypothesis-September-30-2026/paper.pdf}
{the primary preprint}. Such an input can sharpen a prime-summation
error term. It cannot remove the nonzero coefficient in (D4), and
the leading integral in (D2) already follows from the prime number
theorem. No statement of that preprint is used as an additional
axiom in the companion Lean files.

\paragraph{Formalization scope.}
\texttt{OddZeta/HankelForcedDepth.lean} verifies the scaled-coefficient
denominator implication, the terminal model coefficient, and the
special boundary value at $i=p$. The full Lagrange-basis calculation,
its local valuation bounds, and the prime asymptotic above are paper
proofs. They are not represented as a completed Lean proof of
irrationality.
